% !LW recipe = latexmk (xelatex)
\documentclass[12pt,a4paper,UTF8]{ctexart}
\usepackage{SJTUReport}
\usepackage{booktabs,tabularx,longtable,mathtools,enumitem,needspace,tikz}
\usepackage[section]{placeins}
\usetikzlibrary{arrows.meta,positioning,shapes.geometric}
\numberwithin{equation}{section}
\setlength{\emergencystretch}{2em}
\setlength{\headheight}{15pt}
\setlist[enumerate]{itemsep=3pt,topsep=5pt,parsep=0pt}
\setlist[itemize]{itemsep=3pt,topsep=5pt,parsep=0pt}
\renewcommand{\arraystretch}{1.18}
\newcommand{\todoresult}[1]{\textcolor{gray}{#1}}
\newcommand{\ind}{\mathbb I}
\newcommand{\R}{\mathcal R}
\def\subsectionautorefname{节}
\title{\vspace{1.6cm}\heiti\LARGE 基于强化学习终端价值的高速公路\\[0.4em]换电路径与充电功率滚动时域优化\par\vspace{1.2cm}}
\author{\relax}
\date{}

\begin{document}
\cover
\thispagestyle{empty}
\clearpage
\begin{abstract}
面向高速公路长途出行场景，本文研究预约与随机两类用户共同参与的换电站运营问题。车辆在一次行程中可能多次换电，换电站的选择决定沿线各站的服务需求和退回电池状态，充电功率则影响电池恢复服务能力的时间及购电费用。本文对预约用户的换电路径和换电站充电功率进行滚动时域优化，利用强化学习估计预测域的终端价值，综合考虑终端库存和预测域外未完成服务的影响。模型根据车辆续航构建候选网络，以日前换电路径作为在线调整的基准，并结合电池 SOC、等待窗口、服务顺序和站级总功率约束描述运营过程。充电功率与换电路径采用不同的更新周期，每轮执行当前时间段的决策，再根据新观测更新模型。价值函数通过运营轨迹中的实际收益训练，并以分段线性形式纳入优化目标。数值实验分别围绕换电路径与充电功率联合优化、RL 终端价值和运行条件设置对照，为评价所提方法提供统一的实验方案。

\medskip
\noindent\textbf{关键词：}高速公路；换电站；滚动时域优化；充电功率；强化学习；终端价值
\end{abstract}
\clearpage
\tableofcontents
\clearpage

\section{引言}

电动汽车的长途出行需要沿高速公路获得连续、可靠的补能服务。换电站通过交付已经充好的电池缩短车辆的现场补能时间，其服务能力同时取决于车辆到站分布和站内电池库存。长途车辆可能在多个站点换电，前一次选站会改变后续可达站点；每次换电还会向站内退回一块具有特定 SOC 的电池，使车辆的换电安排与各站未来的充电需求相互影响。

运营中同时存在预约用户和随机到站用户。预约信息为需求安排提供依据，车辆实际进入高速公路后的位置、SOC 和预计到站时间又会不断更新。随机用户的到站时间和退回电池状态具有不确定性。与此同时，分时电价、设备功率和站级供电能力共同影响充电安排。因此，运营者需要结合实时信息选择预约用户的换电站点，并协调各站充电功率，以平衡服务收入、购电费用、路径调整和预约履约。

换电路径与充电功率具有直接联系。将车辆安排到某一站点，会改变该站的电池需求、退回 SOC 和后续充电负荷；提高某一站点的充电功率，也可能使原本库存紧张的换电安排获得可行的服务时间。本文将两类决策纳入同一滚动时域优化模型，在每轮预测域内统一评价其运营影响。充电功率按较短周期更新，换电路径按较长周期更新，以适应站内状态变化并保持用户换电安排的稳定。

有限预测域会截断电池库存和预约服务的长期影响。预测结束时仍在充电的电池可能服务于后续需求，已经确定的预约行程也可能在预测域外继续换电。终端价值用于评价这些状态对后续运营的影响。已有研究通过近似价值迭代构造预测控制的终端成本，并分析其与预测长度及价值近似误差的关系\cite{moreno2022terminal}。本文结合高速公路换电场景，将终端库存与预测域外未完成服务共同纳入状态价值函数，通过实际运营轨迹训练，并在每轮优化中计算本轮决策对应的终端价值。

本文的研究内容包括三个方面。首先，利用 SOC 分档候选网络描述预约用户的可达换电路径，建立换电路径与充电功率的联合优化模型。其次，在模型中表达满电交付、退回电池充电、预约优先、等待窗口及路径调整费用，使决策与电池和用户状态的变化相互对应。最后，构造包含终端库存和未完成服务的价值函数，并通过联合优化消融、终端价值消融及敏感性分析评价各组成部分的作用。

图~\ref{fig:framework} 给出整体方法。第 2 节介绍问题与基本假设；第 3 节建立候选网络、日前换电路径和滚动时域优化模型；第 4 节介绍终端价值学习；第 5 节介绍数值实验；第 6 节总结全文。附录给出分段线性函数和相关变量乘积的等价约束。

\begin{figure}[!htbp]
\centering
\begin{tikzpicture}[
  box/.style={draw=blue!45!black,rounded corners=2pt,align=center,
    text width=3.8cm,minimum height=1.1cm,fill=blue!3,font=\small},
  arr/.style={-{Latex[length=2mm]},thick,draw=blue!55!black}]
  \node[box] (obs) {实时状态与预测\\车辆、请求、电池、电价};
  \node[box,right=0.75cm of obs] (opt) {滚动时域优化\\换电路径与充电功率};
  \node[box,right=0.75cm of opt] (act) {当前时间段执行\\服务、充电与状态更新};
  \node[box,below=1.1cm of opt] (value) {RL 终端价值\\终端库存与未完成服务};
  \draw[arr] (obs)--(opt);
  \draw[arr] (opt)--(act);
  \draw[arr] (value)--node[right,font=\small]{价值函数}(opt);
  \draw[arr] (act.south)|-node[pos=0.3,right,font=\small]{实际收益与轨迹}(value.east);
  \draw[arr] (act.north)--++(0,0.6)-|node[pos=0.5,above,font=\small]{下一时刻观测}(obs.north);
\end{tikzpicture}
\caption{换电路径与充电功率滚动时域优化方法}
\label{fig:framework}
\end{figure}

\section{问题描述与基本假设}

\subsection{系统与运营场景}
考虑由多个 O--D 对及沿线换电站组成的高速公路网络。车辆从入口进入后沿下游方向行驶，经必要的换电站补能并从出口离开。各站配备固定数量的电池和充电设备。一次换电取走满电电池并退回带有剩余电量的电池，退回电池经过充电后继续服务其他用户。因此，换电站的选择既影响车辆的后续行程，也影响各站电池 SOC 和充电需求。

系统服务预约用户和随机用户。预约用户提供 O--D、计划入口时间和入口 SOC，研究中的预约用户均具有满足续航条件的完整换电路径。尚未进入高速公路的用户由预约信息描述；进入高速公路后，使用其位置、SOC 和预计到站时间更新换电安排。已经到站且正在等待的请求保留原到站时刻、截止时刻和退回 SOC。随机用户直接到站提出请求，未来随机需求由当前预测信息描述。

\subsection{基本假设}
为描述上述运营过程，作如下假设。
\begin{enumerate}[label=\textbf{假设\arabic*.},leftmargin=4.4em]
  \item \textbf{道路与行驶。}车辆沿高速公路下游方向行驶，并按照给定换电站序完成行程。节点间距离和行驶能耗为已知参数，换电路径满足续航和出口最低 SOC 要求。研究场景中的入口 SOC 与初始路径相容，给定行驶能耗下已确定的换电安排在两次路径更新之间保持续航可行。
  \item \textbf{电池规格与库存。}可交换电池具有统一容量和兼容规格，服务就绪 SOC 归一化为 $1$。每次换电取走和退回各一块电池，站内电池数量保持不变。SOC $1$ 表示运营规定的服务目标。
  \item \textbf{充电过程。}充电效率为已知参数，电池 SOC 按输入电量线性变化。不同站点可以具有不同充电效率，各充电设备和站点分别具有给定功率上限。
  \item \textbf{换电操作。}换电视为瞬时操作。交付电池达到服务就绪 SOC，退回电池进入对应槽位并参加后续充电。槽位记录当时接入的电池状态。
  \item \textbf{信息与预测。}每轮可获得当前电池状态、等待请求及在途车辆的位置和 SOC。车辆到站时间采用当前 ETA 预测，未来随机需求和电价采用决策时刻可得的信息。同一用户沿下游方向的 ETA 按行驶顺序递增，预测信息随滚动过程更新。
  \item \textbf{到站顺序。}同一站点同类请求的原始到站时刻互不相同。同一时间段内可以存在多个请求，其先后次序由原始到站时刻确定。
\end{enumerate}

\subsection{运营规则}
预约请求优先于随机请求，同类请求按原到站时刻先后服务。每个请求自到站时刻起具有给定的最大等待时间，运营者可以在该窗口内选择服务时间，并为其匹配可用满电电池。预约在截止时刻仍未获得服务时，产生一次预约失败成本，该用户后续换电需求随之结束；随机请求在等待截止后流失。已完成的换电和行驶过程按实际记录保留，在线路径调整针对车辆后续的换电安排。换电安排实际发生变化时产生路径调整成本。

\section{模型建立}

将运营时间划分为长度为 $\delta$ 的时间段，记
\begin{equation}
t_n=n\delta,\qquad \mathcal I_n=[t_n,t_{n+1}),\qquad
\mathcal T_\ell=\{\ell,\ldots,\ell+H-1\}.
\label{eq:time_grid}
\end{equation}
其中 $\ell$ 为当前轮，$H$ 为预测时间段数。本轮预测范围为 $[t_\ell,t_{\ell+H})$，在 $t_n$ 处理服务及换电，随后在 $\mathcal I_n$ 内充电。每轮只执行当前时间段，并在下一时刻使用新观测重新优化。默认 $\delta=5$~min，能量计算时取 $\delta=1/12$~h。

充电功率每段更新，换电路径每 $K$ 段允许更新一次。以运营起点为共同基准，定义已知参数
\begin{equation}
g_\ell=\ind\{\ell\bmod K=0\},\qquad K=3.
\label{eq:update_clock}
\end{equation}
因此路径更新周期为 $K\delta=15$~min。$g_\ell=1$ 时联合优化在途用户的路径与充电功率；$g_\ell=0$ 时，整段预测采用当前换电站序并优化充电功率。尚未进入高速公路的用户沿用日前路径，进入后在下一路径更新时刻参与调整。图~\ref{fig:updates} 示意两种决策的更新安排。

\begin{figure}[!htbp]
\centering
\begin{tikzpicture}[x=2.55cm,y=0.7cm,font=\small]
\draw[-{Latex[length=2mm]},thick] (0,0)--(5.15,0);
\foreach \x/\lab in {0/0,1/5,2/10,3/15,4/20,5/25}{
 \draw (\x,-0.08)--(\x,0.08);
 \node[above] at (\x,0.13) {\lab~min};
 \fill[teal!70!black] (\x,-0.7) circle (2.3pt);
}
\foreach \x in {0,3}{\fill[orange!90!black] (\x,-1.5) rectangle ++(0.10,0.16);}
\node[anchor=east] at (-0.14,-0.7) {充电};
\node[anchor=east] at (-0.14,-1.4) {路径};
\node[below,align=center] at (2.5,-1.9) {每轮更新车辆、电池和请求状态，并执行当前时间段决策};
\end{tikzpicture}
\caption{充电功率与换电路径的更新安排}
\label{fig:updates}
\end{figure}

表~\ref{tab:notation} 按集合与索引、参数、决策变量及辅助状态变量、派生量与价值函数列出主要符号。预约用户按 O--D 对分组，以 $\mathcal P$ 表示 O--D 对集合，$p\in\mathcal P$ 为组索引，$k$ 为该组内的用户索引。用户由 $(p,k)$ 唯一标识，编号在滚动过程中保持不变。到站请求统一以 $r$ 标识；某一请求所在站点记为 $i(r)$。费用和价格参数统一采用小写 $c$，合计费用采用大写 $C$，服务收入记为 $I$。

\begin{longtable}{@{}p{4.4cm}p{11.1cm}@{}}
\caption{主要符号}\label{tab:notation}\\
\toprule 符号 & 含义 \\ \midrule
\endfirsthead
\caption[]{主要符号（续）}\\
\toprule 符号 & 含义 \\ \midrule
\endhead
\bottomrule\endfoot
\multicolumn{2}{@{}l@{}}{\textbf{集合与索引}}\\*
$\mathcal S;\ i,j$ & 换电站集合及站点索引 \\
$\mathcal B_i;\ b$ & 站点 $i$ 的电池槽位集合及槽位索引 \\
$\mathcal P;\ p$ & O--D 对集合及索引 \\
$\mathcal K_\ell^p;\ k$ & O--D 对 $p$ 下本轮预约用户集合及组内用户索引 \\
$\mathcal K_\ell^{\mathrm{fut},p},\mathcal K_\ell^{\mathrm{enr},p}$ & O--D 对 $p$ 下本轮尚未进入、已经进入高速公路的预约用户集合 \\
$\R,\R_i;\ r$ & 本轮全部请求集合、站点 $i$ 的请求集合及请求索引 \\
$\R_i^A,\R_i^R$ & 站点 $i$ 的预约请求集合与随机请求集合 \\
$\mathcal V^{p,k},\mathcal A^{p,k}$ & 用户 $(p,k)$ 本轮路径网络的节点集合与可用弧集 \\
$\mathcal T_\ell;\ n,\ell$ & 本轮预测时间段集合；$n$ 为时段索引，$\ell$ 为当前轮索引 \\
\midrule
\multicolumn{2}{@{}l@{}}{\textbf{参数}}\\*
$\delta$ & 时间段长度；默认 5~min，能量计算时以 h 计 \\
$H,K$ & 预测域长度与路径更新间隔，均以时间段数计 \\
$o^{p,k},d^p$ & 用户 $(p,k)$ 本轮路径的起点与出口节点 \\
$v_{i,j}^p$ & O--D 对 $p$ 中节点 $i$ 至 $j$ 的行驶 SOC 消耗 \\
$E_B$ & 电池容量（kWh） \\
$\eta_i$ & 站点 $i$ 的充电效率 \\
$\rho_r$ & 请求 $r$ 对应的退回电池 SOC \\
$S_{i,b}^{\mathrm{obs}}$ & 本轮开始时槽位 $b$ 中电池的观测 SOC \\
$\overline P_{i,b}^{\mathrm{ch}},\overline P_{i,n}^{\mathrm{sta}}$ & 设备充电功率上限与站级总充电功率上限（kW） \\
$t_r^{\mathrm{arr}},t_r^{\mathrm{ddl}}$ & 请求 $r$ 的到站时刻与等待截止时刻；未来到站采用本轮 ETA \\
$\Delta$ & 请求自到站起允许的最大等待时间 \\
$\bar y_{i,j}^{p,k}$ & 用户 $(p,k)$ 的日前换电路径是否使用弧 $(i,j)$ \\
$c_{i,n}^{\mathrm{ch}},c_{i,n}^{\mathrm{sw}}$ & 充电电价与换电服务单价（元/kWh） \\
$c^{\mathrm{adj}},c^{\mathrm{fail}}$ & 单次路径调整成本与单次预约失败成本（元/次） \\
$\beta$ & 终端价值权重，主方法取 1 \\
\midrule
\multicolumn{2}{@{}l@{}}{\textbf{决策变量及辅助状态变量}}\\*
$y_{i,j}^{p,k}$ & 用户 $(p,k)$ 的本轮换电路径是否使用弧 $(i,j)$，二元变量 \\
$P_{i,b,n}$ & 站点 $i$ 的设备 $b$ 在时间段 $n$ 的电网侧输入功率（kW） \\
$\alpha_{b,r,n}$ & 是否在 $t_n$ 用槽位 $b$ 的电池服务请求 $r$，二元变量 \\
$S_{i,b,n}^{-},S_{i,b,n}^{+}$ & 槽位 $b$ 中电池在 $t_n$ 换电前、后的 SOC，由状态递推确定 \\
$\chi_{p,k}$ & 用户 $(p,k)$ 的换电站序在本轮是否改变的二元标记 \\
$f_r$ & 预约请求 $r$ 是否在本预测域内触发用户失败的二元标记 \\
$Q_{r,n}$ & 请求 $r$ 在 $t_n$ 是否已到站、可等待服务的二元标记 \\
$w_r$ & 请求 $r$ 在预测终点是否仍有效且尚未完成的二元标记 \\
\midrule
\multicolumn{2}{@{}l@{}}{\textbf{派生量与价值函数}}\\*
$z_{r,n},Z_r(n),Z_r$ & 请求 $r$ 在 $t_n$ 的服务量、该时刻前的累计服务量及整个预测域内的累计服务量，均由匹配变量汇总 \\
$I$ & 本预测域内的服务收入（元） \\
$C^{\mathrm{ch}},C^{\mathrm{adj}},C^{\mathrm{fail}}$ & 预测域内的充电费用、本轮路径调整费用及预测域内的预约失败费用（元） \\
$s_\ell,s_{\ell+H}$ & 本轮观测运营状态与由本轮决策确定的预测终端状态 \\
$V_\theta(s)$ & 强化学习估计的状态终端价值（元） \\
\end{longtable}

\subsection{基于 SOC 分档的候选网络}
\label{sec:network}
对每个 O--D 对 $p$，以入口、沿线换电站和出口构造网络。车辆路径
$o^p\to i_1\to\cdots\to i_q\to d^p$
表示依次在 $i_1,\ldots,i_q$ 换电，弧所跨越的其他站点不被访问。将入口 SOC 划分为若干区间，分别预先生成候选网络，以减少在线需要处理的路径组合。

节点集合为 $\mathcal V^p=\{o^p\}\cup\mathcal S^p\cup\{d^p\}$，所有节点沿下游方向排列。以 $D_{i,j}^p$ 和 $v_{i,j}^p$ 表示节点间距离和行驶 SOC 消耗，以 $\underline s_d^p$ 表示出口 SOC 下限。最小换电间距 $\underline D$ 用于优先排除距离过近的首段和站间换电安排，出口末段按续航条件处理。对每个 SOC 区间 $h$，按以下四步生成候选网络。
\begin{enumerate}[label=\textbf{步骤\arabic*.},leftmargin=4.2em]
\item \textbf{生成原始弧。}若以该 SOC 区间上限电量能够从入口到达下游站点 $j$，则连接入口与 $j$；若 $v_{i,j}^p\le1$，则连接下游站点 $i,j$；若 $v_{i,d^p}^p+\underline s_d^p\le1$，则连接站点与出口。入口可直接到达出口时加入相应弧。
\item \textbf{生成完整路径。}在原始网络中枚举所有入口至出口的完整路径。
\item \textbf{剪枝过近换电弧。}找出距离小于 $\underline D$ 的首段弧和站间弧，按距离从小到大考察。若当前仍存在至少一条不包含该弧的完整路径，则删除该弧；否则保留。
\item \textbf{形成候选网络。}合并最终保留路径中的弧，得到 SOC 区间对应的弧集 $\mathcal A^{h,p}$。
\end{enumerate}

在线使用时，根据用户实际入口 SOC 进一步检查首段续航。对在途用户，删除已经经过的节点，以实时位置建立虚拟起点，按实时 SOC 检查其到下游站点的首段弧。最小换电间距始终以入口或最近一次实际换电位置为基准。对已经在站等待的用户，将当前站作为后续路径起点，后续行程以该次换电完成后的满电状态计算；当前等待请求在第~\ref{sec:model}~节中保留，并与后续服务关联。由此得到用户弧集 $\mathcal A^{p,k}$、起点 $o^{p,k}$ 以及对应出口 $d^p$。

\subsection{日前换电路径计算}
\label{sec:dayahead}
日前换电路径根据预约用户的 O--D、入口 SOC 和候选网络计算。先使用实际预约入口 SOC 检查首段弧，并保留满足出口 SOC 要求的完整路径。设所得完整路径集合为 $\mathcal Y^{p,k}$，其元素由弧指示量 $y$ 表示，则日前路径取为
\begin{equation}
\bar y^{p,k}\in\arg\min_{y\in\mathcal Y^{p,k}}
\sum_{(j,i)\in\mathcal A^{p,k}:\,i\in\mathcal S}y_{j,i}.
\label{eq:dayahead_path}
\end{equation}
该目标选择换电次数最少的完整路径。当多条路径的换电次数相同时，依次比较其第一、第二及后续换电站的位置，优先选择更靠近出口的站点；站点位置相同的记录使用固定网络顺序区分。得到的 $\bar y^{p,k}$ 及对应换电站序列作为在线路径调整的基准。该计算仅需要车辆与道路信息，输入用户均具有满足上述条件的完整可达路径。

\subsection{滚动时域优化模型}
\label{sec:model}

\paragraph{请求与输入信息。}
对每个 $p\in\mathcal P$，令 $\mathcal K_\ell^p=\mathcal K_\ell^{\mathrm{fut},p}\cup\mathcal K_\ell^{\mathrm{enr},p}$。对于用户 $k\in\mathcal K_\ell^p$，对每条以换电站 $i$ 为终点的可用弧 $(j,i)\in\mathcal A^{p,k}$，建立记录 $r=(p,k,j,i)$，其站点为 $i(r)=i$，到站时刻由用户在该站的 ETA 给定，退回 SOC 为
\begin{equation}
\rho_r=
\begin{cases}
s_{\mathrm{dep}}^{p,k}-v_{o^{p,k},i}^{p,k},&j=o^{p,k},\ \text{用户正在行驶或尚未进入高速},\\
1-v_{j,i}^{p},&j\text{ 为换电站}.
\end{cases}
\label{eq:return_soc}
\end{equation}
其中 $v_{o^{p,k},i}^{p,k}$ 为本轮起点到站点 $i$ 的已知 SOC 消耗，$s_{\mathrm{dep}}^{p,k}$ 为本轮起点的已知车辆 SOC，未到用户使用预约信息及入口前观测，在途用户使用实时值。用户当前在站等待时，后续首段按第二种情况计算。各记录的 $\rho_r$ 在求解前已知，并由可达弧筛选保证 $0\le\rho_r<1$。

真实等待请求使用独立记录，保留观测到站时间、原截止时间和退回 SOC，预约等待记录同时保留所属用户 $(p,k)$；当前站在路径网络中作为起点，因此同一次正在等待的换电只出现一条记录。随机请求包括当前真实等待请求和预测在 $[t_\ell,t_{\ell+H})$ 到站的请求。令 $\R_i^A,\R_i^R$ 分别为站点 $i$ 的预约和随机请求记录，$\R_i=\R_i^A\cup\R_i^R$，并以 $\R^A,\R^R,\R$ 表示各站相应集合的并集。预约集合覆盖本轮路径上全部可能的后续换电，包括 ETA 位于预测域外的记录；域外随机需求以第 4 节的预测特征提供。

以 $\R_{p,k}^A$ 表示用户 $(p,k)$ 的全部预约记录。对于普通路径记录，定义选择量 $a_r=y_{j,i}^{p,k}$；真实等待请求和随机请求取 $a_r=1$。同一用户同一站点对应的不同入弧是互斥的路径选择。等待截止时刻及服务窗口参数为
\begin{equation}
t_r^{\mathrm{ddl}}=t_r^{\mathrm{arr}}+\Delta,
\qquad e_{r,n}=\ind\{t_r^{\mathrm{arr}}\le t_n\le t_r^{\mathrm{ddl}}\}.
\label{eq:service_window}
\end{equation}
当前时刻 $e_{r,\ell}$ 由真实到站和等待记录确定，只有实际在站且仍在等待窗口内的请求取 $1$。以后时段采用式~\eqref{eq:service_window} 的预测值。新预测请求的 ETA 根据当前观测更新，真实等待请求保持原始时刻。

\paragraph{决策变量。}
主要变量为路径 $y_{i,j}^{p,k}\in\{0,1\}$ 和充电功率 $P_{i,b,n}\ge0$。以 $\alpha_{b,r,n}\in\{0,1\}$ 表示在 $t_n$ 用站点 $i(r)$ 的槽位 $b$ 服务请求 $r$，其中 $b\in\mathcal B_{i(r)}$。换电前后 SOC 分别为 $S_{i,b,n}^{-},S_{i,b,n}^{+}$。另设路径改变标记 $\chi_{p,k}$、预约失败标记 $f_r$、排队标记 $Q_{r,n}$ 和终端未完成服务标记 $w_r$，均为二元变量。服务量及其累计值定义为
\begin{equation}
z_{r,n}=\sum_{b\in\mathcal B_{i(r)}}\alpha_{b,r,n},\qquad
Z_r(n)=\sum_{\substack{m\in\mathcal T_\ell\\m<n}}z_{r,m},\qquad
Z_r=\sum_{n\in\mathcal T_\ell}z_{r,n}.
\label{eq:service_sums}
\end{equation}
这些量由匹配变量直接计算。满足下述容量约束时，$z_{r,n},Z_r(n),Z_r$ 均取 $0$ 或 $1$。

\paragraph{目标函数。}
本轮目标为
\begin{equation}
\max\quad I-C^{\mathrm{ch}}-C^{\mathrm{adj}}-C^{\mathrm{fail}}
+\beta V_\theta(s_{\ell+H}),
\label{eq:objective}
\end{equation}
其中各项分别为
\begin{subequations}\label{eq:objective_terms}
\begin{align}
I&=E_B\sum_{n\in\mathcal T_\ell}\sum_{r\in\R}
c_{i(r),n}^{\mathrm{sw}}(1-\rho_r)z_{r,n},\label{eq:income}\\
C^{\mathrm{ch}}&=\delta\sum_{n\in\mathcal T_\ell}\sum_{i\in\mathcal S}
c_{i,n}^{\mathrm{ch}}\sum_{b\in\mathcal B_i}P_{i,b,n},\label{eq:charging_cost}\\
C^{\mathrm{adj}}&=c^{\mathrm{adj}}\sum_{p\in\mathcal P}\sum_{k\in\mathcal K_\ell^{\mathrm{enr},p}}\chi_{p,k},\label{eq:adjustment_cost}\\
C^{\mathrm{fail}}&=c^{\mathrm{fail}}\sum_{r\in\R^A}f_r.\label{eq:failure_cost}
\end{align}
\end{subequations}
单位服务收入按交付电量 $E_B(1-\rho_r)$ 计算，充电费用按输入电量计算，路径调整和预约失败分别以用户的一次调整和一次失败计费。终端价值由第 4 节的模型评价终端库存及预测域外未完成服务，默认 $\beta=1$。各收入与费用按实际发生时刻所在时间段归属。

\paragraph{（1）路径和续航约束。}
对每个用户，路径满足
\begin{equation}
\sum_{j:(i,j)\in\mathcal A^{p,k}}y_{i,j}^{p,k}
-\sum_{j:(j,i)\in\mathcal A^{p,k}}y_{j,i}^{p,k}
=\begin{cases}
1,&i=o^{p,k},\\-1,&i=d^p,\\0,&\text{其他节点},
\end{cases}
\quad p\in\mathcal P,\ k\in\mathcal K_\ell^p.
\label{eq:flow}
\end{equation}
式~\eqref{eq:flow} 对 $i\in\mathcal V^{p,k}$ 成立，其中 $\mathcal V^{p,k}$ 为 $\mathcal A^{p,k}$ 的节点集合，当前用户满足 $o^{p,k}\ne d^p$。全文空集合求和取零。弧沿下游方向排列，式~\eqref{eq:flow} 因而确定一条起点至出口的路径。可用弧按第~\ref{sec:network}~节检查首段 SOC、站间满电续航和出口 SOC。构造 $\mathcal A^{p,k}$ 时，对能直接驶向出口并满足出口 SOC 要求的节点，仅保留其通向出口的出弧。当前站等待请求作为必须处理的记录保留，其服务完成条件通过式~\eqref{eq:predecessor} 连接后续换电。

令 $x_i^{p,k}=\sum_{j:(j,i)\in\mathcal A^{p,k}}y_{j,i}^{p,k}$ 表示后续是否访问站点 $i$。$\bar x_i^{p,k}$ 为日前站序对应量，$x_i^{p,k,\mathrm{pub}}$ 为当前已经向用户确定的站序对应量；两者都先剔除已经完成的换电，当前正在等待的站点由其独立请求保留。比较集合取当前与参照站序所涉及站点的并集，缺少的访问量补零。对各 $p\in\mathcal P$，按用户状态和更新时间施加
\begin{equation}
\begin{aligned}
x_i^{p,k}&=\bar x_i^{p,k},
&&k\in\mathcal K_\ell^{\mathrm{fut},p},\\
|x_i^{p,k}-x_i^{p,k,\mathrm{pub}}|&\le g_\ell,
&&k\in\mathcal K_\ell^{\mathrm{enr},p}.
\end{aligned}
\label{eq:path_freeze}
\end{equation}
单向道路上的访问站点决定站序，因此该关系保持实际换电安排，并允许起点随车辆位置更新。对在途用户定义差异量 $D_i^{p,k}=|x_i^{p,k}-x_i^{p,k,\mathrm{pub}}|$，其精确线性表达及路径改变标记为
\begin{equation}
\begin{aligned}
D_i^{p,k}&\ge x_i^{p,k}-x_i^{p,k,\mathrm{pub}},&
D_i^{p,k}&\ge x_i^{p,k,\mathrm{pub}}-x_i^{p,k},\\
D_i^{p,k}&\le x_i^{p,k}+x_i^{p,k,\mathrm{pub}},&
D_i^{p,k}&\le 2-x_i^{p,k}-x_i^{p,k,\mathrm{pub}},\\
\chi_{p,k}&\ge D_i^{p,k}\quad(\text{各 }i),&
\chi_{p,k}&\le\sum_iD_i^{p,k},\qquad 0\le D_i^{p,k}\le1.
\end{aligned}
\label{eq:path_change}
\end{equation}
因此 $\chi_{p,k}$ 精确表示站序是否变化。每轮仅评价本次可执行的路径调整。

\paragraph{（2）服务容量和前后关系。}
每个请求至多获得一次服务，每个槽位在同一时刻至多交付一次：
\begin{equation}
Z_r\le a_r\quad(r\in\R),\qquad
\sum_{r\in\R_i}\alpha_{b,r,n}\le1
\quad(i\in\mathcal S,b\in\mathcal B_i,n\in\mathcal T_\ell).
\label{eq:assignment_capacity}
\end{equation}
对普通预约记录 $r=(p,k,j,i)$，以前一站 $j$ 的服务作为其后续服务条件。令 $\mathcal P_r$ 为用户 $(p,k)$ 在站点 $j$ 的所有可能换电记录；当前正在该站等待时，$\mathcal P_r$ 只包含该真实等待记录。随机请求、真实等待请求，以及非在站等待用户从本轮起点到达的首次请求，均取 $\mathcal P_r=\varnothing$。已在此前完成的服务由当前起点状态继承。定义
\begin{equation}
B_{r,n}=\begin{cases}
1,&\mathcal P_r=\varnothing,\\
\displaystyle\sum_{q\in\mathcal P_r}Z_q(n),&\mathcal P_r\ne\varnothing,
\end{cases}
\qquad z_{r,n}\le B_{r,n}.
\label{eq:predecessor}
\end{equation}
路径和容量约束使 $B_{r,n}$ 为二元量，后一站服务只能在前一站完成后的时间段发生。本轮 ETA 给出预测到站窗口，式~\eqref{eq:predecessor} 表达服务的先后关系；车辆行驶和等待对 ETA 的后续影响由下一轮观测更新。预测窗口是给定 ETA 下的运营近似。

\paragraph{（3）预约失败与随机流失。}
设 $h_r=\ind\{t_r^{\mathrm{ddl}}<t_{\ell+H}\}$。真实等待请求在本轮开始时仍有效，所有记录的截止时刻均不早于 $t_\ell$。对预约请求，定义此前截止的失败数量
\begin{equation}
F_r^{\mathrm{pre}}=
\sum_{\substack{q\in\R_{p,k}^A\\t_q^{\mathrm{ddl}}<t_r^{\mathrm{ddl}}}}f_q,
\qquad
F_{p,k}=\sum_{r\in\R_{p,k}^A}f_r\le1.
\label{eq:failure_order}
\end{equation}
其中 $(p,k)$ 为请求 $r$ 对应用户。$h_r=0$ 时取 $f_r=0$；$h_r=1$ 时施加
\begin{equation}
\begin{aligned}
f_r&\le a_r,& f_r&\le1-Z_r,& f_r&\le1-F_r^{\mathrm{pre}},\\
f_r&\ge a_r-Z_r-F_r^{\mathrm{pre}}.
\end{aligned}
\label{eq:failure_exact}
\end{equation}
同一用户沿所选路径的 ETA 和截止时刻依次递增，以上关系使第一个超过窗口且未获服务的请求产生一次失败，后续请求随该用户退出服务过程。截止恰好等于预测终点时，由下一轮处理。随机请求的流失量直接定义为 $h_r(1-Z_r)$。

\paragraph{（4）等待和服务顺序。}
为按每个时刻实际可参与分配的请求确定优先顺序，定义
\begin{equation}
F_{p,k,n}=\sum_{\substack{q\in\R_{p,k}^A\\t_q^{\mathrm{ddl}}<t_n}}f_q,
\qquad
F_{r,n}=\begin{cases}F_{p,k,n},&r\in\R_{p,k}^A,\\0,&r\in\R^R.\end{cases}
\label{eq:failure_before_time}
\end{equation}
恰在 $t_n$ 截止的请求先参加当次服务，因此此处使用严格小于。排队标记 $Q_{r,n}$ 满足
\begin{equation}
\begin{aligned}
Q_{r,n}&\le a_r,& Q_{r,n}&\le e_{r,n},&
Q_{r,n}&\le1-Z_r(n),\\
Q_{r,n}&\le B_{r,n},& Q_{r,n}&\le1-F_{r,n},\\
Q_{r,n}&\ge a_r+e_{r,n}+B_{r,n}-Z_r(n)-F_{r,n}-2,\\
z_{r,n}&\le Q_{r,n}.
\end{aligned}
\label{eq:queue}
\end{equation}
这些约束精确表示所选请求已经到站、尚未服务、前一次必要换电已经完成且用户尚未失败。$Q_{r,n}$ 使用该时刻之前的失败结果。

在同一站点定义 $q\prec r$：当 $q$ 为预约而 $r$ 为随机请求，或者二者同类且 $t_q^{\mathrm{arr}}<t_r^{\mathrm{arr}}$ 时成立。服务顺序满足
\begin{equation}
z_{r,n}\le 1-Q_{q,n}+z_{q,n},
\qquad q\prec r,\quad q,r\in\R_i,\quad n\in\mathcal T_\ell.
\label{eq:priority}
\end{equation}
较高优先级请求仍在等待时，较低优先级请求的服务以较高优先级请求同刻获得服务为条件。运营者可以在窗口内选择服务时点；同刻有多个服务时，由匹配变量决定各自使用的电池。

\paragraph{（5）满电交付和充电。}
槽位初始状态取观测 SOC。充电与换电满足
\begin{subequations}\label{eq:battery_dynamics}
\begin{align}
S_{i,b,\ell}^{-}&=S_{i,b}^{\mathrm{obs}},\label{eq:initial_soc}\\
S_{i,b,n}^{-}&\ge\alpha_{b,r,n},
&&r\in\R_i,\label{eq:full_delivery}\\
S_{i,b,n}^{+}&=S_{i,b,n}^{-}
-\sum_{r\in\R_i}(1-\rho_r)\alpha_{b,r,n},\label{eq:swap_soc}\\
S_{i,b,n+1}^{-}&=S_{i,b,n}^{+}
+\frac{\eta_i\delta}{E_B}P_{i,b,n},\label{eq:charge_soc}\\
0\le S_{i,b,n}^{-},S_{i,b,n}^{+}&\le1,
\qquad 0\le S_{i,b,\ell+H}^{-}\le1,\label{eq:soc_bounds}\\
0\le P_{i,b,n}&\le\overline P_{i,b}^{\mathrm{ch}},\label{eq:device_power}\\
\sum_{b\in\mathcal B_i}P_{i,b,n}&\le\overline P_{i,n}^{\mathrm{sta}}.
\label{eq:station_power}
\end{align}
\end{subequations}
除已标注的请求索引外，上述关系对所有 $i\in\mathcal S,b\in\mathcal B_i,n\in\mathcal T_\ell$ 成立。式~\eqref{eq:full_delivery} 与 SOC 上界保证交付前电池达到 $1$，式~\eqref{eq:swap_soc} 将退回 SOC 写入相应槽位。功率在当前时间段内保持恒定，SOC 上界限制本段输入电量。每个请求可以选择任意满足条件的槽位，模型同时确定其后续充电功率。

\paragraph{（6）终端库存与未完成服务。}
预测终端库存取 $S_{i,b,\ell+H}^{-}$，即终点服务之前的状态。对预约请求，未完成服务标记定义为
\begin{equation}
\begin{aligned}
w_r&\le a_r,&w_r&\le1-Z_r,&w_r&\le1-F_{p,k},\\
w_r&\ge a_r-Z_r-F_{p,k},
&&r\in\R_{p,k}^A.
\end{aligned}
\label{eq:terminal_pending}
\end{equation}
因此，所选路径中尚未服务且用户仍有效的请求均被保留，包括前一站尚未服务的后续预约及预计在预测域外到站的预约。对随机请求，取
\begin{equation}
w_r=(1-h_r)(1-Z_r),\qquad r\in\R^R.
\label{eq:terminal_random}
\end{equation}
在式~\eqref{eq:queue} 中排队标记的上下界取 $n=\ell+H$，利用该时刻以前的累计服务和失败量，得到终端等待标记 $Q_{r,\ell+H}$；此时只计算状态。$Q_{r,\ell+H}=1$ 的预约属于 $w_r=1$ 的未完成服务，状态编码中以同一条记录及其等待标记表示。

以 $t^{\mathrm{end}}$ 表示新增需求停止的时刻，已给定预约及其全部后续请求仍在集合中保留。为统一预测域内和预测终点的运营结束条件，将未完成服务标记扩展到各边界 $n\in\mathcal T_\ell\cup\{\ell+H\}$：
\begin{equation}
w_r(n)=\begin{cases}
a_r[1-Z_r(n)][1-F_{p,k,n}],&r\in\R_{p,k}^A,\\
\ind\{t_r^{\mathrm{ddl}}\ge t_n\}[1-Z_r(n)],&r\in\R^R.
\end{cases}
\label{eq:pending_each_boundary}
\end{equation}
其中 $w_r=w_r(\ell+H)$，上述二元乘积在终点与式~\eqref{eq:terminal_pending}--\eqref{eq:terminal_random} 一致。定义各边界的结束标记及终端结束标记
\begin{equation}
d_n^{\mathrm{end}}=
\ind\{t_n\ge t^{\mathrm{end}}\}
\prod_{r\in\R}[1-w_r(n)],\qquad
d^{\mathrm{end}}=d_{\ell+H}^{\mathrm{end}}.
\label{eq:operating_end}
\end{equation}
新增需求停止且全部服务完成或失败后，运营在首个满足 $d_n^{\mathrm{end}}=1$ 的决策边界结束。最后一次服务所在时间段完成既定状态转移，后续时间段的功率满足
\begin{equation}
P_{i,b,n}\le\overline P_{i,b}^{\mathrm{ch}}(1-d_n^{\mathrm{end}}),
\qquad i\in\mathcal S,\ b\in\mathcal B_i,\ n\in\mathcal T_\ell.
\label{eq:power_at_end}
\end{equation}
结束后的服务收入与各项费用均为零。上述二元标记和乘积关系按附录线性化。最后一次换电完成后驶向出口的用户已经完成全部服务。终端状态由电池 SOC、上述未完成请求及等待标记、换电安排、时间、电价和后续需求预测共同构成，其具体编码在第 4 节给出。

式~\eqref{eq:flow}--\eqref{eq:power_at_end}、变量范围和第 4 节价值函数共同构成优化问题。主要约束为线性关系，价值计算及其中有界变量乘积按附录表示后，得到混合整数线性优化模型。

\paragraph{滚动执行与记账。}
在 $t_\ell$ 读取真实在站请求、电池状态、车辆位置和最新预测，建立本轮模型。路径更新时刻向在途用户确定新的换电安排，并据此记录调整成本；其他时刻保持已确定的站序。执行模型选择的当前服务和匹配后，在 $\mathcal I_\ell$ 内使用当前充电功率。区间内到站记录在下一格点进入服务安排，原到站时刻用于计算等待时间。截止位于 $\mathcal I_\ell$ 内且尚未服务的真实请求在该段结算超时；恰在下一格点截止的请求留到下一轮先参加服务。

下一轮电池状态由实际服务和充电更新；真实等待请求保留原标识及截止时间；车辆位置和 ETA 由新观测更新。未来预测记录在实际到站后以对应观测替换，状态中始终区分已观测请求与预测请求。若上游预约实际失败，该用户后续服务一起结束。新的当前状态成为下一轮初始条件，新的路径更新参数按式~\eqref{eq:update_clock} 计算。

\section{终端价值学习}
\label{sec:value}

强化学习用于估计预测域的终端价值，重点考虑终端库存和预测域外未完成服务。相同的电池总电量可能对应不同的满电供应能力；相同的库存也可能面对不同的预约到站时间、等待窗口和后续换电需求。价值函数将这些信息共同映射为一个终端评价，并通过实际运营轨迹中的累计净收益进行训练。

\subsection{状态与终端价值}
\label{sec:value_state}

以 $s(\tau)$ 表示时刻 $\tau$ 的完整运营状态，包括电池 SOC、真实等待请求、未完成预约及其换电安排，以及时间、电价和需求预测。已进入高速公路的预约还具有位置、SOC 和 ETA 等观测。对于预测终端，使用式~\eqref{eq:battery_dynamics} 的库存、式~\eqref{eq:terminal_pending}--\eqref{eq:terminal_random} 的未完成请求以及终点等待标记构造相同含义的状态。下文的用户和请求集合均取该状态时刻的集合，$\mathcal K^p$ 表示该时刻 O--D 对 $p$ 下的预约用户集合，省略其时间标记。

为适应不同数量的请求，采用共享分段线性函数和求和生成固定维度表示。对任意状态时刻 $\tau$，令 $\nu_r(\tau)$ 汇集请求站点和前一站的位置、用户 O--D、退回 SOC、相对到站与截止时间，以及该请求是否正在等待。预约记录还包含该用户当前换电站序和可得车辆信息；在预测终端采用本轮确定的站序及给定车辆预测。终端需要继续完成的站点用
\begin{equation}
x_i^{p,k,\mathrm{end}}=\sum_{r\in\R_{p,k}^A:\,i(r)=i}w_r
\label{eq:terminal_visits}
\end{equation}
表示，其中保留仍在等待的当前站点。站序按下游方向排列，已完成的站点随 $w_r=0$ 移除；请求等待特征取 $Q_{r,\ell+H}$。由此，终端站序和等待状态直接由本轮决策确定，其余车辆预测与时刻特征为给定输入。用户编号仅用于将同一用户的请求归组。

以 $w_r(\tau)$ 表示该状态下请求仍有效且尚未完成，以 $Q_r(\tau)$ 表示其中已经进入等待的请求。真实状态取观测和现有预约，预测终端分别取 $w_r$ 和 $Q_{r,\ell+H}$。令 $\xi(\tau)$ 包含运营时间、距离需求期结束的时间、当前与后续电价、随机需求预测和路径更新相位。路径更新相位由时间索引模 $K$ 得到；终端以 $\ell+H$ 的相位计算。状态编码为
\begin{equation}
\begin{aligned}
f_B(\tau)&=\mathop{\mathrm{concat}}_{i\in\mathcal S}
\left(\sum_{b\in\mathcal B_i}\phi_B(S_{i,b}^{-}(\tau))\right),\\
h_{p,k}(\tau)&=\sum_{r\in\R_{p,k}^A}w_r(\tau)\psi_A\bigl(\nu_r(\tau)\bigr),\\
f_A(\tau)&=\sum_{p\in\mathcal P}\sum_{k\in\mathcal K^p}
\left[\phi_A(h_{p,k}(\tau))-\phi_A(0)\right],\\
f_R(\tau)&=\sum_{r\in\R^R}w_r(\tau)\psi_R\bigl(\nu_r(\tau)\bigr),\\
\Phi_\theta(s(\tau))&=\bigl[f_B(\tau),f_A(\tau),f_R(\tau),\xi(\tau)\bigr].
\end{aligned}
\label{eq:state_encoding}
\end{equation}
其中 $\phi_B,\phi_A,\psi_A,\psi_R$ 为固定维度的分段线性映射，使用仿射变换和 ReLU 构造。站内电池使用共享映射，以反映 SOC 分布；预约先按用户汇总再整体编码，以保留同一行程多次换电的关联。$\phi_A(h_{p,k})-\phi_A(0)$ 使已经没有后续服务的用户贡献为零。等待状态作为请求自身的一个特征，预约等待请求在 $f_A$ 中记录一次。域外随机需求作为 $\xi$ 的预测输入。

以 $s_\ell=s(t_\ell)$、$s_{\ell+H}=s(t_{\ell+H})$ 分别表示当前状态及预测终端状态。真实与预测状态采用相同特征定义和归一化参数，归一化参数由训练集确定后固定。式~\eqref{eq:state_encoding} 将完整状态映射为固定维度的网络输入，其表示能力通过第 5 节的终端价值消融实验评价。

原始网络输出及实际使用的价值函数定义为
\begin{equation}
\begin{aligned}
\widetilde V_\theta(s)
&=d_0+\boldsymbol v_0^\top\Phi_\theta(s)
+\sum_{m=1}^{M}a_m\max\{0,\boldsymbol v_m^\top\Phi_\theta(s)+d_m\},\\
V_\theta(s)&=(1-d^{\mathrm{end}}(s))\widetilde V_\theta(s).
\end{aligned}
\label{eq:value_network}
\end{equation}
结束标记在当前状态由实际服务记录确定，在预测终端由式~\eqref{eq:operating_end} 确定。价值函数的训练参数 $\theta$ 包含网络和状态映射中的参数。每轮优化时固定这些参数，终端库存和未完成请求则随本轮决策变化。状态映射中的有界变量乘积、ReLU 及结束标记乘积按附录展开后，终端价值与业务约束共同进入同一个优化问题。

\subsection{实际收益与状态更新}
\label{sec:reward}
每轮由优化模型返回当前路径、功率及相应服务方案，首段执行结果决定实际收益。令 $\mathcal H_\ell^{\mathrm{sw}}$ 为在 $t_\ell$ 实际获得服务的请求，$\mathcal H_\ell^{\mathrm{fail}}$ 为截止时刻位于 $\mathcal I_\ell$ 内且实际超时的预约请求，$N_\ell^{\mathrm{adj}}$ 为本轮实际改变换电站序的用户数，则
\begin{equation}
\begin{aligned}
r_\ell={}&E_B\sum_{r\in\mathcal H_\ell^{\mathrm{sw}}}
c_{i(r),\ell}^{\mathrm{sw}}(1-\rho_r)\\
&-\delta\sum_{i\in\mathcal S}c_{i,\ell}^{\mathrm{ch}}
\sum_{b\in\mathcal B_i}P_{i,b,\ell}^{*}
-c^{\mathrm{adj}}N_\ell^{\mathrm{adj}}
-c^{\mathrm{fail}}|\mathcal H_\ell^{\mathrm{fail}}|.
\end{aligned}
\label{eq:actual_reward}
\end{equation}
其中 $\rho_r$ 为实际退回 SOC，当前功率为执行值。实际在站请求参与首段服务，预测中的未来服务和失败在其真实发生后进入对应时间段的收益。跨边界等待请求在实际服务时记录一次收入；失败预约的后续请求结束，失败成本只记录一次。

设 $a_\ell^*$ 汇集本轮实际执行的路径更新、功率和服务安排，以 $\varepsilon_{\ell+1}$ 表示该时间段到站记录、车辆观测及下一轮预测的更新，则
\begin{equation}
s_{\ell+1}=\mathcal F(s_\ell,a_\ell^*,\varepsilon_{\ell+1}).
\label{eq:state_transition}
\end{equation}
映射 $\mathcal F$ 按第~\ref{sec:model}~节的换电、充电、等待及结算规则更新完整运营记录；价值评价时再按式~\eqref{eq:state_encoding} 编码。预测中尚未真实发生的到站以新观测校正，因此训练样本来自实际执行的状态和收益。求解器返回的匹配直接用于该时间段执行，状态转移保留相应电池变化。

每个运营场景包含规定的新增需求期。需求期结束后继续处理已给定预约和等待请求，直至全部服务完成或用户失败。后续服务所需的电价和 ETA 信息延伸至该场景结束。剩余电池残值取零。若场景在第 $L$ 个状态结束，采用未折扣回报
\begin{equation}
G_\ell=\sum_{n=\ell}^{L-1}r_n,\qquad G_L=0.
\label{eq:return_target}
\end{equation}
该回报与式~\eqref{eq:objective} 的金额单位和结算项目一致，预测域内收益与终端之后的价值按时间划分。

\subsection{价值函数训练}
\label{sec:training}
采用交替采样与价值拟合的训练过程。首先使用零终端价值的滚动时域优化生成完整运营轨迹。随后，以各状态对应的回报 $G_\ell$ 训练价值网络，损失函数为
\begin{equation}
\mathcal L(\theta)=\frac{1}{|\mathcal D|}
\sum_{(s_\ell,G_\ell)\in\mathcal D}
\bigl(V_\theta(s_\ell)-G_\ell\bigr)^2.
\label{eq:value_loss}
\end{equation}
其中 $\mathcal D$ 保存当前采样批次的完整状态记录与回报，训练时按当前参数 $\theta$ 重新计算 $\Phi_\theta(s_\ell)$，以同时更新状态映射和价值网络。网络更新后，将新价值函数用于下一批场景中的滚动时域优化，再采集轨迹并拟合。这样，优化模型根据当前终端价值生成决策，价值学习根据这些决策的实际结果更新评价。

\begin{table}[!htbp]
\centering
\caption{终端价值学习流程}
\label{tab:training_algorithm}
\begin{tabularx}{\linewidth}{@{}lX@{}}
\toprule 步骤 & 内容 \\ \midrule
1 & 固定训练、验证和测试场景划分，初始化用于优化的终端价值为零。\\
2 & 在一个采样批次内固定价值函数及特征提取参数，按 5~min 与 15~min 的更新安排执行滚动时域优化，记录完整运营轨迹。\\
3 & 由式~\eqref{eq:actual_reward} 计算各段实际收益，按式~\eqref{eq:return_target} 从场景结束向前累计回报。\\
4 & 最小化式~\eqref{eq:value_loss}，更新状态映射和价值网络；归一化统计在训练阶段确定并固定。\\
5 & 用更新后的价值函数重复采样和拟合，按验证场景的实际净收益选择网络。\\
6 & 固定选定网络，在独立测试场景上评价运营表现和求解时间。\\
\bottomrule
\end{tabularx}
\end{table}

该过程可视为由滚动时域优化生成策略、由回报拟合进行评价的近似策略迭代。每次求解期间价值函数保持固定，参数更新发生在采样批次之间。训练场景覆盖不同需求强度、SOC 分布和等待情况，网络结构、批次规模和训练预算在实验设置中记录。价值近似与滚动优化之间的相互作用通过验证场景的实际运营表现评价。

\section{数值实验}
\label{exp:section}

数值实验从换电路径与充电功率的联合优化、RL 终端价值和运行条件三个方面评价所提方法。各组实验使用相同的道路网络、预约用户集合、初始换电路径和初始电池状态。训练集、验证集与测试集按运营场景划分，各组对照在相同测试场景上运行。测试阶段固定价值网络参数，需求预测仅使用决策时刻可获得的信息。

计算环境见表~\ref{exp:environment}，主要实验参数见表~\ref{exp:parameters}。每个时间段为 5~min，充电功率每段更新，换电路径每 15~min 允许更新一次。采用学习终端价值的主方法取 $\beta=1$。预测域长度、场景规模、机器配置和其余参数按实际实验设置填写。

\begin{table}[!htbp]
  \centering
  \caption{计算环境}
  \label{exp:environment}
  \small
  \begin{tabularx}{0.88\linewidth}{@{}lX@{}}
    \toprule
    项目 & 配置 \\
    \midrule
    机器型号 & 待填 \\
    CPU 型号 & 待填 \\
    GPU 型号及显存 & 待填 \\
    内存容量 & 待填 \\
    操作系统 & 待填 \\
    Python 版本 & 待填 \\
    Gurobi 版本 & 待填 \\
    PyTorch 版本 & 待填 \\
    \bottomrule
  \end{tabularx}
\end{table}

\begin{table}[!htbp]
  \centering
  \caption{主要实验参数}
  \label{exp:parameters}
  \small
  \setlength{\tabcolsep}{5pt}
  \begin{tabularx}{\linewidth}{@{}Xll@{}}
    \toprule
    参数 & 符号或单位 & 设置 \\
    \midrule
    站点数量与位置 & 个，km & 待填 \\
    各站电池数量与初始 SOC & $|\mathcal B_i|$，无量纲 & 待填 \\
    电池容量 & $E_B$，kWh & 待填 \\
    各站充电效率 & $\eta_i$ & 待填 \\
    设备功率上限 & $\overline P_{i,b}^{\mathrm{ch}}$，kW & 待填 \\
    站级总功率上限 & $\overline P_{i,n}^{\mathrm{sta}}$，kW & 待填 \\
    用户等待时间 & min & 待填 \\
    预约规模、O--D 及入口 SOC 分布 & 用户数及分布参数 & 待填 \\
    随机用户需求与预测设置 & 到站率及预测参数 & 待填 \\
    充电电价与换电服务单价 & $c_{i,n}^{\mathrm{ch}},c_{i,n}^{\mathrm{sw}}$，元/kWh & 待填 \\
    路径调整与预约失败成本 & $c^{\mathrm{adj}},c^{\mathrm{fail}}$，元/次 & 待填 \\
    时间段长度与充电更新周期 & $\delta$ & 5~min \\
    路径更新周期 & $K\delta$ & 15~min \\
    预测域长度 & $H$，段 & 待填 \\
    主方法终端价值权重 & $\beta$ & 1 \\
    需求期长度 & h & 待填 \\
    训练、验证及测试场景数量 & 个 & 待填 \\
    独立训练次数与测试重复次数 & 次 & 待填 \\
    训练预算 & 运营轨迹数及更新次数 & 待填 \\
    单次求解时限与相对最优间隙 & s，\% & 待填 \\
    \bottomrule
  \end{tabularx}
\end{table}

各方法均在需求期结束后继续处理已有预约和等待请求，直至完成或失败，并采用相同的终止状态和结算方式。运营净收益为实际服务收入减去充电、路径调整和预约失败成本。预约失败率按未完成预约行程的用户数占预约用户总数计算；同一用户的预约失败只计一次。随机用户服务率为获得服务的随机请求数占实际随机请求总数的比例。平均等待时间按已获服务请求的实际服务时刻与原始到站时刻之差计算，并同时报告服务率以反映未获服务用户的情况。路径调整次数记录实际更新给用户的换电站序列变化。各项运营指标报告重复实验的均值与标准差，求解时间报告单轮平均值和第 95 百分位数。样本分母为零的指标记为“不适用”。

\FloatBarrier
\subsection{换电路径与充电功率联合优化消融实验}
\label{exp:joint}

为考察两项决策的作用，设置表~\ref{exp:joint_methods}所列四组方法。各组终端价值均取零，使用相同预测域长度、需求输入和服务规则。服务时间与满电电池匹配均由优化模型确定。采用日前路径的组别在运营过程中沿用日前计算的换电站序列；优化换电路径的组别按规定的更新周期调整在途预约用户的换电安排。

\begin{table}[!htbp]
  \centering
  \caption{换电路径与充电功率联合优化的对照设置}
  \label{exp:joint_methods}
  \small
  \begin{tabularx}{\linewidth}{@{}Xlll@{}}
    \toprule
    方法 & 换电路径 & 充电功率 & 终端价值 \\
    \midrule
    基准方法 & 日前路径 & 基准充电规则 & 0 \\
    仅优化换电路径 & 在线优化 & 基准充电规则 & 0 \\
    仅优化充电功率 & 日前路径 & 在线优化 & 0 \\
    联合优化 & 在线优化 & 在线优化 & 0 \\
    \bottomrule
  \end{tabularx}
\end{table}

基准充电规则将站级总功率额度按设备数量等分，实际功率同时满足设备上限和电池可充电量限制：
\begin{equation}
 P_{i,b,n}^{\mathrm{base}}=
 \min\left\{
 \overline P_{i,b}^{\mathrm{ch}},\;
 \frac{\overline P_{i,n}^{\mathrm{sta}}}{|\mathcal B_i|},\;
 \frac{E_B(1-S_{i,b,n}^{+})}{\eta_i\delta}
 \right\}.
 \label{eq:baseline_power}
\end{equation}
该规则使用本时刻完成换电后的 SOC。每台设备保留等分得到的功率额度，满电电池对应设备的功率为零。采用基准规则的组别在预测域内施加 $P_{i,b,n}=(1-d_n^{\mathrm{end}})P_{i,b,n}^{\mathrm{base}}$，使各运营时间段使用基准功率，结束后的时间段功率为零。其分段线性表达见附录。式中 $\delta$ 以小时计，因而第三项的单位为 kW。

表~\ref{exp:joint_results}记录运营指标，表~\ref{exp:joint_costs}记录收入与成本分解。通过基准方法与两组单项优化的比较，考察路径调整和充电安排各自的作用；通过两组单项优化与联合优化的比较，考察同时优化两项决策的运营表现。各项比较使用相同测试场景的配对结果。

\begin{table}[!htbp]
  \centering
  \caption{联合优化消融实验的运营结果}
  \label{exp:joint_results}
  \small
  \setlength{\tabcolsep}{3pt}
  \begin{tabularx}{\linewidth}{@{}Xcccccc@{}}
    \toprule
    方法 & \shortstack{净收益\\元} & \shortstack{预约\\失败率/\%} & \shortstack{随机用户\\服务率/\%} & \shortstack{等待时间\\min} & \shortstack{调整次数\\次} & \shortstack{求解时间\\均值/P95，s} \\
    \midrule
    基准方法 & 待填 & 待填 & 待填 & 待填 & 待填 & 待填 \\
    仅优化换电路径 & 待填 & 待填 & 待填 & 待填 & 待填 & 待填 \\
    仅优化充电功率 & 待填 & 待填 & 待填 & 待填 & 待填 & 待填 \\
    联合优化 & 待填 & 待填 & 待填 & 待填 & 待填 & 待填 \\
    \bottomrule
  \end{tabularx}
\end{table}

\begin{table}[!htbp]
  \centering
  \caption{联合优化消融实验的收入与成本分解（元）}
  \label{exp:joint_costs}
  \small
  \begin{tabularx}{\linewidth}{@{}Xcccc@{}}
    \toprule
    方法 & 服务收入 & 充电成本 & 路径调整成本 & 预约失败成本 \\
    \midrule
    基准方法 & 待填 & 待填 & 待填 & 待填 \\
    仅优化换电路径 & 待填 & 待填 & 待填 & 待填 \\
    仅优化充电功率 & 待填 & 待填 & 待填 & 待填 \\
    联合优化 & 待填 & 待填 & 待填 & 待填 \\
    \bottomrule
  \end{tabularx}
\end{table}

\todoresult{待填：根据运营指标和成本分解，分析换电路径优化、充电功率优化及二者联合优化的作用，并给出重复实验的统计结果。}

\FloatBarrier
\subsection{RL终端价值消融实验}
\label{exp:terminal}

本节各组均对换电路径与充电功率进行联合优化，采用相同预测长度、域内需求信息、状态观测和业务约束。对照方法仅改变预测终点的价值表达，如表~\ref{exp:value_methods}所示。终端库存反映各站的后续供给能力，未完成服务包括预测结束时仍在等待的请求，以及预约用户预计在预测域外发生的换电需求。

\begin{table}[!htbp]
  \centering
  \caption{终端价值的对照设置}
  \label{exp:value_methods}
  \small
  \begin{tabularx}{\linewidth}{@{}p{2.6cm}XX@{}}
    \toprule
    方法 & 终端价值 & 比较内容 \\
    \midrule
    零终端价值 & $V=0$ & 终端价值的总体作用 \\
    简单库存价值 & $c^{\mathrm{inv}}E_B\sum_{i,b}S_{i,b}^{-}$ & 人工库存估值与学习价值的差别 \\
    RL 库存价值 & 使用库存及公共外生信息 & 库存价值学习的作用 \\
    完整 RL 终端价值 & 使用库存、未完成服务及公共外生信息 & 未完成服务信息的作用 \\
    \bottomrule
  \end{tabularx}
\end{table}

表中 SOC 均取预测终点服务发生前的值。简单库存价值的系数 $c^{\mathrm{inv}}$ 的单位为元/kWh，根据训练集或验证集确定后固定。各组终端函数在运营终止状态均取零。两组 RL 方法使用同一网络结构形式，公共外生信息包括时间、电价、后续随机需求预测和路径更新相位；RL 库存价值以终端库存及这些公共信息作为输入，完整 RL 终端价值进一步使用未完成服务及相应的用户换电安排。输入维度按各自特征确定，输入归一化参数由训练数据计算后固定。

两个学习方案分别生成训练轨迹并更新价值网络，使用相同的训练场景划分、运营轨迹预算和参数更新预算。网络的最终选择依据验证集表现，测试阶段固定网络参数，并在相同测试场景上评价。表~\ref{exp:training}记录训练设置，表~\ref{exp:value_results}记录运营结果。图~\ref{exp:value_figure}预留训练过程和运营表现的展示位置。

\begin{table}[!htbp]
  \centering
  \caption{终端价值网络的训练设置}
  \label{exp:training}
  \small
  \begin{tabularx}{\linewidth}{@{}Xcc@{}}
    \toprule
    项目 & RL 库存价值 & 完整 RL 终端价值 \\
    \midrule
    输入维度 & 待填 & 待填 \\
    隐藏单元数 & 待填 & 待填 \\
    参数总数 & 待填 & 待填 \\
    训练场景数量 & 待填 & 待填 \\
    采样批次数与每批轨迹数 & 待填 & 待填 \\
    每批参数更新次数 & 待填 & 待填 \\
    回归批量大小 & 待填 & 待填 \\
    优化器与学习率 & 待填 & 待填 \\
    独立训练次数 & 待填 & 待填 \\
    验证频率与网络选择规则 & 待填 & 待填 \\
    实际训练耗时 & 待填 & 待填 \\
    \bottomrule
  \end{tabularx}
\end{table}

\begin{table}[!htbp]
  \centering
  \caption{RL 终端价值消融实验的运营结果}
  \label{exp:value_results}
  \small
  \setlength{\tabcolsep}{3pt}
  \begin{tabularx}{\linewidth}{@{}Xcccccc@{}}
    \toprule
    方法 & \shortstack{净收益\\元} & \shortstack{预约\\失败率/\%} & \shortstack{随机用户\\服务率/\%} & \shortstack{等待时间\\min} & \shortstack{调整次数\\次} & \shortstack{求解时间\\均值/P95，s} \\
    \midrule
    零终端价值 & 待填 & 待填 & 待填 & 待填 & 待填 & 待填 \\
    简单库存价值 & 待填 & 待填 & 待填 & 待填 & 待填 & 待填 \\
    RL 库存价值 & 待填 & 待填 & 待填 & 待填 & 待填 & 待填 \\
    完整 RL 终端价值 & 待填 & 待填 & 待填 & 待填 & 待填 & 待填 \\
    \bottomrule
  \end{tabularx}
\end{table}

\begin{figure}[!htbp]
  \centering
  \fbox{\begin{minipage}[c][32mm][c]{0.9\linewidth}
    \centering
    待填：价值网络训练过程及终端价值消融结果\\[4pt]
    横轴与单位、各方法曲线及重复实验的不确定性范围待填。
  \end{minipage}}
  \caption{终端价值学习过程及运营表现（待填）}
  \label{exp:value_figure}
\end{figure}

\todoresult{待填：比较各类终端价值对应的运营指标，分析终端库存和预测域外未完成服务的作用，并结合求解时间评价计算代价。}

\FloatBarrier
\subsection{敏感性分析}
\label{exp:sensitivity}

敏感性分析考察预测域长度、需求强度、电池库存配置和站级总功率上限。每组实验改变一个因素，其余设置保持基准值；充电与路径更新周期始终为 5~min 和 15~min。表~\ref{exp:sensitivity_settings}记录各因素的基准值和变化范围。预测域长度以 $H$ 个时间段表示，并同时报告对应的小时数 $H\delta$。

\begin{table}[!htbp]
  \centering
  \caption{敏感性分析的参数设置}
  \label{exp:sensitivity_settings}
  \small
  \begin{tabularx}{\linewidth}{@{}Xlll@{}}
    \toprule
    因素 & 基准值 & 变化范围 & 取值数量 \\
    \midrule
    预测域长度 $H$ & 待填 & 待填 & 待填 \\
    预约与随机需求相对基准的倍数 & 待填 & 待填 & 待填 \\
    各站电池数量 $|\mathcal B_i|$ & 待填 & 待填 & 待填 \\
    站级总功率上限相对基准的倍数 & 待填 & 待填 & 待填 \\
    \bottomrule
  \end{tabularx}
\end{table}

预测域长度实验比较零终端价值与完整 RL 终端价值，对各个 $H$ 使用同一组测试场景，记录运营表现和求解时间。该比较用于分析预测域截断与终端价值之间的关系。需求实验按同一倍数调整预约规模和随机到站强度，保持两类用户的比例及其他分布参数一致。库存实验改变各站电池数量，并保持初始 SOC 分布和站级功率上限一致。功率实验按同一倍数调整各站总功率上限，保持电池数量和设备功率上限一致。

完整 RL 终端价值在需求强度、电池配置或功率配置发生变化时，使用相应配置下的训练与验证场景训练，训练预算保持一致；每组测试固定对应网络。预测域长度比较使用同一配置下固定的价值网络。所有测试场景均与训练及验证场景分离。表~\ref{exp:sensitivity_results}提供结果记录格式，每个参数取值及每种被比较的方法分别填列一行，图~\ref{exp:sensitivity_figure}展示主要指标随参数的变化。

\begin{table}[!htbp]
  \centering
  \caption{敏感性分析结果（各参数取值分别填列）}
  \label{exp:sensitivity_results}
  \small
  \setlength{\tabcolsep}{3pt}
  \begin{tabularx}{\linewidth}{@{}Xcccccc@{}}
    \toprule
    因素、取值及方法 & \shortstack{净收益\\元} & \shortstack{预约\\失败率/\%} & \shortstack{随机用户\\服务率/\%} & \shortstack{等待时间\\min} & \shortstack{调整次数\\次} & \shortstack{求解时间\\均值/P95，s} \\
    \midrule
    预测域长度：待填 & 待填 & 待填 & 待填 & 待填 & 待填 & 待填 \\
    需求强度：待填 & 待填 & 待填 & 待填 & 待填 & 待填 & 待填 \\
    电池库存配置：待填 & 待填 & 待填 & 待填 & 待填 & 待填 & 待填 \\
    站级总功率：待填 & 待填 & 待填 & 待填 & 待填 & 待填 & 待填 \\
    \bottomrule
  \end{tabularx}
\end{table}

\begin{figure}[!htbp]
  \centering
  \fbox{\begin{minipage}[c][36mm][c]{0.9\linewidth}
    \centering
    待填：预测域长度、需求强度、电池库存和站级总功率的敏感性结果\\[4pt]
    各参数取值、运营指标曲线及求解时间曲线待填。
  \end{minipage}}
  \caption{不同运行条件下的运营表现与求解时间（待填）}
  \label{exp:sensitivity_figure}
\end{figure}

\todoresult{待填：分别分析预测域长度、需求强度、电池数量和站级总功率变化时的收益、服务表现及求解时间，并说明各项结论所对应的参数范围。}

\section{结论}

本文围绕高速公路换电站网络中的预约与随机用户，建立换电路径与充电功率的滚动时域优化方法。候选网络根据 SOC 和道路条件描述车辆可达的换电安排，日前路径提供在线调整的基准。优化模型将路径选择、充电、电池状态及用户服务条件联系起来，以不同周期更新换电路径和充电功率，并按实际服务、充电、路径改变和预约失败记录运营收益。

强化学习用于估计预测域的终端价值，将终端库存与预测域外未完成服务共同纳入评价。模型保留电池 SOC 分布、等待信息及预约后续换电需求，并通过分段线性网络表达其价值。价值函数由完整运营轨迹中的实际回报训练，网络参数在每轮优化中固定。数值实验按联合优化消融、RL 终端价值消融和敏感性分析组织，具体结果与定量结论将在实验完成后填写。

本文以每轮 ETA 预测描述到站窗口，以瞬时换电和给定充电效率表达站内运行。状态编码和价值函数具有有限维度，其效果需要结合实际运营指标和求解时间评价。后续研究可进一步结合交通过程、换电作业时间及电池特性，分析更丰富的运行条件下路径与充电的协调方法。


\appendix
\section{线性化方法}
\label{lin:appendix}

本附录给出正文中分段线性函数、变量乘积及逻辑关系的等价约束。每轮求解时，网络参数和各表达式的有效上下界均为已知常数，新增变量仅用于表达相应关系。

\paragraph{二元变量与有界连续变量的乘积。}
设 $x\in[L,U]$、$b\in\{0,1\}$，引入连续变量 $w=bx$，其等价表达为
\begin{equation}
\begin{aligned}
Lb\le w&\le Ub,\\
x-U(1-b)\le w&\le x-L(1-b).
\end{aligned}
\label{lin:product}
\end{equation}
当 $b=0$ 时得到 $w=0$，当 $b=1$ 时得到 $w=x$。该表达适用于含负值的区间，可用于请求状态标记与连续特征的乘积。若 $d\in\{0,1\}$ 表示运营已经结束，原始网络输出为 $\widetilde V\in[L_V,U_V]$，则终端价值 $V=(1-d)\widetilde V$ 也按式~\eqref{lin:product} 表达，其中取 $b=1-d$。因此，运营结束状态的终端价值严格等于零。

\paragraph{ReLU 与价值网络。}
对 $h=\max\{0,x\}$，若 $L<0<U$，引入二元变量 $\xi$，有
\begin{equation}
\begin{aligned}
h&\ge 0, & h&\ge x,\\
h&\le U\xi, & h&\le x-L(1-\xi),
\qquad \xi\in\{0,1\}.
\end{aligned}
\label{lin:relu}
\end{equation}
当 $U\le0$ 时直接置 $h=0$；当 $L\ge0$ 时直接置 $h=x$。对式~\eqref{eq:value_network} 中每个仿射输入 $x_m=\boldsymbol v_m^\top\Phi_\theta(s)+d_m$ 分别使用上述约束，并将对应的输出变量 $h_m$ 代入
\begin{equation}
\widetilde V=d_0+\boldsymbol v_0^\top\Phi_\theta(s)+\sum_m a_m h_m.
\label{lin:network}
\end{equation}
上述等式和约束共同确定网络值，适用于任意符号的输出权重 $a_m$。状态编码中的分段线性特征按相同方法展开；其输出作为后续网络的输入，从而将特征提取与价值计算一并纳入优化模型。

\paragraph{有限个二元条件的逻辑关系。}
设 $b_j\in\{0,1\}$，$j\in\mathcal J$，其中 $\mathcal J$ 为有限非空集合。二元变量 $q$ 表示至少一个条件成立，二元变量 $a$ 表示全部条件成立，其等价约束分别为
\begin{equation}
\begin{aligned}
q&\ge b_j\quad(j\in\mathcal J), & q&\le\sum_{j\in\mathcal J}b_j,\\
a&\le b_j\quad(j\in\mathcal J), & a&\ge\sum_{j\in\mathcal J}b_j-|\mathcal J|+1.
\end{aligned}
\label{lin:logic}
\end{equation}
空集合对应的逻辑或取 $0$，逻辑与取 $1$。正文涉及路径改变、请求有效性及运营结束的条件可按其逻辑定义逐项使用这些约束。

\paragraph{基准充电规则。}
对式~\eqref{eq:baseline_power}，记已知常数与仿射表达式为
\begin{equation}
\begin{aligned}
\bar p_{i,b,n}
&=\min\left\{\overline P_{i,b}^{\mathrm{ch}},
\frac{\overline P_{i,n}^{\mathrm{sta}}}{|\mathcal B_i|}\right\},\\
q_{i,b,n}&=\frac{E_B(1-S_{i,b,n}^{+})}{\eta_i\delta},
\qquad Q_i=\frac{E_B}{\eta_i\delta}.
\end{aligned}
\label{lin:baseline_bounds}
\end{equation}
由 $0\le S_{i,b,n}^{+}\le1$ 可得 $0\le q_{i,b,n}\le Q_i$。若 $\bar p_{i,b,n}=0$，直接置 $P_{i,b,n}^{\mathrm{base}}=0$；若 $\bar p_{i,b,n}\ge Q_i$，直接置 $P_{i,b,n}^{\mathrm{base}}=q_{i,b,n}$。在 $0<\bar p_{i,b,n}<Q_i$ 时，引入 $\zeta_{i,b,n}\in\{0,1\}$，写成
\begin{equation}
\begin{aligned}
P_{i,b,n}^{\mathrm{base}}&\le\bar p_{i,b,n},\\
P_{i,b,n}^{\mathrm{base}}&\le q_{i,b,n},\\
P_{i,b,n}^{\mathrm{base}}&\ge\bar p_{i,b,n}\zeta_{i,b,n},\\
P_{i,b,n}^{\mathrm{base}}&\ge q_{i,b,n}-(Q_i-\bar p_{i,b,n})\zeta_{i,b,n}.
\end{aligned}
\label{lin:baseline_min}
\end{equation}
当 $\zeta_{i,b,n}=1$ 时功率取设备与站级额度共同确定的上限；当 $\zeta_{i,b,n}=0$ 时功率取本段可充电量对应的值。两种情况均满足两项上界，因而精确表示二者的最小值。运营结束条件与基准功率的乘积 $P_{i,b,n}=(1-d_n^{\mathrm{end}})P_{i,b,n}^{\mathrm{base}}$ 按式~\eqref{lin:product} 表达，其中基准功率的有效区间为 $[0,\min\{\bar p_{i,b,n},Q_i\}]$。

\paragraph{有效上下界。}
SOC 的边界取 $[0,1]$，请求状态标记取 $[0,1]$，功率使用正文中的设备和站级上限。请求数量由当前有限请求集合限定；到站时间、截止时间和等待时间特征根据本轮已知时刻及有效等待窗口确定范围。对于含状态标记的特征，先求连续部分的边界，再使用式~\eqref{lin:product}。若某层输入向量为 $\boldsymbol e$，$x=\boldsymbol v^\top\boldsymbol e+d$ 且 $e_j\in[\underline e_j,\overline e_j]$，令 $v_j^+=\max\{v_j,0\}$、$v_j^-=\min\{v_j,0\}$，则可取
\begin{equation}
\begin{aligned}
L&=d+\sum_j\left(v_j^+\underline e_j+v_j^-\overline e_j\right),\\
U&=d+\sum_j\left(v_j^+\overline e_j+v_j^-\underline e_j\right).
\end{aligned}
\label{lin:interval}
\end{equation}
ReLU 输出的区间为 $[\max\{0,L\},\max\{0,U\}]$。依次传递各层区间，可得到隐藏单元和最终输出的有限上下界，并据此固定本轮线性化常数。若求解前能够由业务约束得到更紧的边界，则使用相应边界以缩小连续松弛范围。


\begin{thebibliography}{9}
\bibitem{moreno2022terminal}
Moreno-Mora F, Beckenbach L, Streif S.
Predictive Control with Learning-Based Terminal Costs Using Approximate Value Iteration[EB/OL].
arXiv:2212.00361, 2022.
\url{https://arxiv.org/abs/2212.00361}.
\end{thebibliography}
\end{document}
